\documentclass[12pt]{article}
\usepackage[T1]{fontenc}
\usepackage[utf8]{inputenc}
\usepackage{lmodern}
\usepackage{textcomp}
\usepackage{enumitem}
\usepackage{geometry}
\geometry{margin=1in}
\begin{document}
\begin{center}
Answer Key - Chemistry - Undergraduate Level Assessment 9 \
\small (Answers are ordered exactly as in the question paper.)
\end{center}

\section*{Multiple Choice}
\begin{enumerate}[label=\arabic*.]
\item \textbf{Answer:} C
\\ \textit{Options:} A) 101.1 x $10^7$ T-1 s-1; B) 26.75 x $10^7$ T-1 s-1; C) 7.081 x $10^7$ T-1 s-1; D) 0.9377 x $10^7$ T-1 s-1
\\ \textit{Note:} The magnetogyric ratio is calculated by dividing the NMR frequency (198.52 MHz) by the magnetic field strength (750 MHz), yielding a value of approximately 7.081 x $10^7$ T^-1 s^-1.
\item \textbf{Answer:} D
\\ \textit{Options:} A) T$_{1}$, unlike T$_{2}$, is sensitive to very low-frequency molecular motions.; B) T$_{2}$, unlike T$_{1}$, is sensitive to very low-frequency molecular motions.; C) T$_{1}$, unlike T$_{2}$, is sensitive to molecular motions at the Larmor frequency.; D) T$_{2}$, unlike T$_{1}$, is sensitive to molecular motions at the Larmor frequency.
\\ \textit{Note:} The correct answer is based on the fact that T$_{1}$ (spin-lattice relaxation time) is primarily influenced by energy transfer to the surrounding lattice, while T$_{2}$ (spin-spin relaxation time) is affected by interactions between nuclear spins and molecular motions at the Larmor frequency, which can lead to faster relaxation and potentially shorter T$_{2}$ values compared to T$_{1}$.
\item \textbf{Answer:} D
\\ \textit{Options:} A) I only; B) II only; C) III only; D) I and III only
\\ \textit{Note:} The normal modes of a carbon dioxide molecule that are infrared-active must involve a change in the dipole moment, which occurs in the asymmetric stretching and bending modes, but not in the symmetric stretching mode, as it does not produce a net change in the dipole moment.
\item \textbf{Answer:} A
\\ \textit{Options:} A) Neutrons; B) Alpha particles; C) Beta particles; D) X-rays
\\ \textit{Note:} Cobalt-60 is produced through a nuclear reaction where cobalt-59 absorbs a neutron, leading to its transformation into cobalt-60, which is utilized in radiation therapy for cancer treatment.
\item \textbf{Answer:} C
\\ \textit{Options:} A) Nuclear magnetic resonance; B) Infrared; C) Raman; D) Ultraviolet-visible
\\ \textit{Note:} Raman spectroscopy is a light-scattering technique because it involves the inelastic scattering of monochromatic light, typically from a laser, which results in a shift in energy corresponding to vibrational modes of molecules.
\end{enumerate}

\section*{True / False}
\begin{enumerate}[label=\arabic*.]
\setcounter{enumi}{5}
\item \textbf{Answer:} False
\\ \textit{Options:} A) True; B) False
\\ \textit{Note:} The statement is false because the hydrogen atom in CHCl 3 undergoes rapid intermolecular exchange due to hydrogen bonding, leading to averaged chemical environments that typically result in a single peak in its spectrum rather than distinct peaks.
\item \textbf{Answer:} True
\\ \textit{Options:} A) True; B) False
\\ \textit{Note:} At standard temperature and pressure, bromine (Br$_{2}$) is a diatomic molecule that exhibits a red-brown color and is indeed a volatile liquid due to its relatively low boiling point of 58.8°C.
\item \textbf{Answer:} True
\\ \textit{Options:} A) True; B) False
\\ \textit{Note:} The statement is true because the 55 Mn nuclide, being a stable isotope with specific nuclear configurations, has three distinct energy levels corresponding to its nuclear spin states that can be populated under certain conditions.
\item \textbf{Answer:} False
\\ \textit{Options:} A) True; B) False
\\ \textit{Note:} The statement is false because gadolinium (Gd) has the electron configuration [Xe] 4 $f^7$ 5 $d^1$, indicating that it does possess a partially filled 4 f subshell with seven electrons.
\item \textbf{Answer:} False
\\ \textit{Options:} A) True; B) False
\\ \textit{Note:} The equilibrium polarization of 13 C nuclei in a 20.0 T magnetic field at 300 K is actually lower than 10.8 x 10^-5 due to the much smaller magnetic moment of 13 C compared to protons, resulting in a significantly reduced polarization.
\end{enumerate}

\section*{Long Form Response}
\begin{enumerate}[label=\arabic*.]
\setcounter{enumi}{10}
\item \textbf{Answer:} The stability of the +1 and +3 oxidation states in group 13 elements varies, with thallium (Tl) exhibiting a notable preference for the +1 state. This preference is largely attributed to relativistic effects, which play a significant role in heavier elements like thallium. In the +1 oxidation state, thallium achieves a stable electron configuration similar to that of the noble gas xenon, while in the +3 state, it experiences increased ionization energy and lower stability due to the higher oxidation state involving the loss of more electrons. Additionally, the +1 state is more stabilized by the inert pair effect, where the s-electrons are held more tightly and are less likely to participate in bonding, thus reinforcing the stability of the +1 oxidation state over the +3 state in thallium.
\\ \textit{Note:} correct because it identifies that thallium's preference for the +1 oxidation state is influenced by relativistic effects, resulting in a stable electron configuration akin to xenon, while the +3 state is less stable due to higher ionization energy and increased oxidation challenges.
\item \textbf{Answer:} In the reaction 3 Cl-(aq) + 4 CrO $4^2$-(aq) + 23 H+(aq) → 3 HClO 2(aq) + 4 Cr$_{3}$+(aq) + 10 H$_{2}$O(l), Cl-(aq) acts as a reducing agent by donating electrons to the chromium species. As Cl- is oxidized to HClO 2, it loses electrons, which facilitates the reduction of CrO $4^2$- to Cr$_{3}$+. This electron transfer is a defining characteristic of reducing agents, as they enable other substances to gain electrons while undergoing oxidation themselves. Therefore, the role of Cl-(aq) as a reducing agent is crucial in driving the redox chemistry of the reaction forward.
\\ \textit{Note:} correct because it accurately describes how Cl-(aq) donates electrons during the reaction, thus undergoing oxidation to HClO 2 while facilitating the reduction of CrO $4^2$- to Cr$_{3}$+, which exemplifies the defining characteristic of a reducing agent.
\end{enumerate}

\vfill
\noindent\textit{End of Answer Key}
\end{document}